\section{讨论与局限}\label{sec:discussion}

这些结果与一个简单的解释相一致。位置恰当的早期决策改变后续轮继承的状态。FRO-2 在严格意义上更保守：首轮翻转更少，其中错误翻转少得多，留给后续轮的真实错误更多。更轻、更干净的首轮是良性方向；更大的残差是被推迟的工作，而不是损害。因此「尽早减少错误」不是一个充分的目标；位移的符号、时点与幅度共同决定结果。

规范方的反对意见值得公平陈述。生产 BIKE-Flip 在每一轮都维持裕度，而我们的干预只触及一轮；两种主张针对不同调度。规范标定的是整条调度，而我们的破坏性情形是对逐轮标定调度施加的中途位移。FRO-2 不是规范自带的首轮保护；它是普通 BF 译码器上的一次性探针。

随机化对照通过把位置从调度中移除来补全论证。对照运行在 BF30 上，这是一个固定阈值 BF 变体，并在每一轮为阈值加入掷币式抖动。对称抖动把失败率从 4.29\% 推到 19.67\%；仅正向的单侧配方落在 50.5\% 与 51.8\% 之间。同一配方用于普通 BF 译码器则不留下可分辨差异（Fisher $p\ge0.54$）。无结构混合不能替代有定位的确定性调度。

本研究有四个边界。全部主要结论依托约 $10^{-4}$ 到 $10^{-3}$ 的失败率，因此不外推生产目标。两个工作点承载主要推断，第三个为补充；其它块尺寸只有行为筛查。偏移网格与被筛调度是我们的构造，跨实现检查是同项目而非独立复现。诊断标签使用真实错误支持集，部署者并不知道它。在生产参数下，抬高重量的窗口比生产重量 134 高 28\%；生产重量证据是决策余量与收敛成本，零失败上界为 $9.4\times10^{-5}$。捕获未被证明；交叠是压力构造。未来工作应把测量推向生产尾部，并在更大密钥池上测试独立实现。

\section{结论}\label{sec:conclusion}

普通 BF 译码器的首轮是一扇陡峭且不对称的双向窗口。保守偏移有益，而宽松侧在两个计数内崩塌。决定条件失败风险的是偏移的轮次与持续方式，而不是单独的幅度。窗口与不对称的假翻转抑制相一致。干预剔除错误决策，而不是加速纠错。它无需重调即跨工作点与重量迁移，在所测捕获倾向构造处止步，并在生产块尺寸上于错误重量抬升后复现；在生产重量本身只有决策余量与收敛代价可测。一个首轮代理模型未能预测另一套调度的效应，随机化也不能复现有定位决策的收益。实践上，逐变体配对实例、以密钥为聚类单元做推断，并在失败可测处标定阈值。
